\section{Reuse Notes and Open Questions}

\paragraph{Why this paper matters for our work.}
This paper is useful less because it solves two-layer networks completely, and more because it replaces parameter optimization by measure evolution. That shift is conceptually important for any architecture where the right large-scale object is not a vector of parameters but a distribution over architectural components.

\paragraph{Potential reuse directions.}
\begin{itemize}[leftmargin=1.5em]
\item The empirical-measure viewpoint may be relevant if we later want to study neural architectures whose effective state is a population over nodes, paths, or motifs.
\item The ``good initialization + PDE limit + PDE convergence'' proof pattern is a reusable template.
\item The Hessian-based stability criteria in Theorems~6--7 suggest a way to reason about stable versus unstable architectural regimes without solving the full finite system.
\end{itemize}

\paragraph{Open questions left by the paper.}
\begin{itemize}[leftmargin=1.5em]
\item Can the dependence of the convergence time on dimension and tolerance be made explicit in generic settings?
\item How much of the clean noisy-DD theory survives without regularization?
\item Can analogous PDE limits be developed for architectures with non-exchangeable geometric structure, where the parameter distribution is not permutation-symmetric?
\end{itemize}

\paragraph{Bottom line.}
The main reusable lesson is not ``SGD always works'' but rather: in some high-width regimes, the right object to analyze is a probability measure evolving under a gradient-flow PDE, and that object may have a much cleaner landscape than the original parameter system.
